\section{Introduction}

A property that completes a rival's monopoly is not an ordinary item in a bilateral exchange. Selling it may finance the seller while giving the buyer a source of rent. Refusing may preserve a blocker while leaving both parties exposed to a third player. If either later becomes insolvent, the deed can pass to a creditor without another negotiation. An economic agent must therefore consider the exchange, the development it enables and the positions that persist after it.

Monopoly makes these dependencies observable over a complete game. Players buy, auction and trade properties, build houses, mortgage collateral, pay rents and liquidate assets. In the implementation studied here, a player-creditor inherits a bankrupt debtor's properties with their buildings intact. The recipient of an estate can consequently matter to players who were not parties to the debt. This paper asks how LLM agents bargain over such dependencies, and what happens between an acquisition, an investment request and a realized economic outcome.

We study twenty complete four-player games involving four configured systems, abbreviated Sol, Gemini, Opus and Grok. The corpus includes all 526 negotiation threads, 209 group-completion episodes and sixty eliminations. Legal menus, response attempts, delivered prompts and canonical events let us relate the terms of an exchange to the opportunities and actions that follow. The unit of inference is the configured system operating in these games, including its response handling and information. It is not an isolated language model.

The results follow an economic sequence. First, agents bargain over what an exchange enables. They discuss a rival's building funds, the value of withholding a completing deed and their own reserves after a purchase. In sixty-one of seventy-four rejections of offers that would complete the recipient's group, both textual readings identify a stated consequence for the rival. Counteroffers commonly preserve the proposed exchange while demanding more cash. The response census places these explanations in context: negotiation stage and recipient configuration account for substantial differences in rejection and countering.

Second, development is the step that converts a collection of deeds into a rent-producing position. The four systems complete similar numbers of groups by trade, but build on sixteen, seventeen, seven and three in their first management windows. The trajectories show mortgage-financed investment, cash-reserve decisions and delayed development, together with requests that do not execute. Groups developed in that first window collect \$53,810 of the \$57,378 subsequently received by trade-completed groups. Following the financing and execution sequence explains how these positions differ after apparently similar acquisitions.

Third, creditors turn another player's failure into an acquisition. Fifty-six of sixty elimination debts exceed even a generous own-asset financing bound. Eventual winners receive 202 of 227 deeds transferred through estates before the final elimination, and estates complete fifty-two groups while play continues. This makes creditor identity a strategic consideration before bankruptcy: participants bargain over keeping a rival solvent, preserving a blocker and preventing another player from inheriting a complete group.

Eight contextual incidents accompany these results. They show a seller changing its ask with the buyer's position, refusal of one's own completion, a public bargaining limit below a private contingency, an omitted financing route, an appeal to withhold a rival's building cash, and withdrawal followed by lower bids after a seller's deterioration. Correction followed by acceptance and a contradictory bankruptcy response provide contrasts. Each transcript follows the statement through its recipient's response and the later economic trajectory. Public deception was permitted by the task; its occurrence and economic context are distinct questions from instruction adherence.

Prior work studies sustained economic operation, state-grounded communication, Monopoly behavior and complete model-plus-execution configurations~\cite{backlund2025vendingbench,li2026arena,wu2026agentmonopoly,hu2026finboardbench,zhang2026dece}. Our contribution is to connect bargaining by all four players to subsequent financing, development, income and creditor succession over complete games. Legal menus and historical feasibility checks establish the opportunities behind those trajectories. The discussion relates the findings to strategic reasoning, opponent modeling, negotiation, planning, adaptation, risk management, collective behavior and evaluation validity, distinguishing observed behavior from questions requiring a different study design.
